\documentclass[11pt]{article}
\usepackage[margin=1in]{geometry}
\usepackage[spanish,es-noshorthands]{babel}
\usepackage[T1]{fontenc}
\usepackage[utf8]{inputenc}
\usepackage{amsmath,amssymb}
\usepackage{tikz}
\usetikzlibrary{shapes.geometric}
\usepackage{multicol}
\usepackage{verbatim}

\setlength{\parindent}{0pt}
\setlength{\parskip}{0.45\baselineskip}

\title{Reducción de imágenes astronómicas: calibración y overscan}
\author{Notas de clase}
\date{}

\begin{document}

\begin{center}
  {\Large\bfseries Reducción de imágenes astronómicas}\par
  {\normalsize Calibración de CCD y corrección de \texttt{overscan}}\par
\end{center}

\section*{1. El detector CCD e imágenes de calibración}

La lectura de un detector CCD puede ser lenta: el sistema debe leer cada píxel y
convertir su señal en un valor digital. Además de la imagen científica, se toman
imágenes de calibración para medir y corregir la respuesta instrumental. Las
imágenes \texttt{bias}, \texttt{dark} y \texttt{flat} describen efectos
diferentes y no deben confundirse.

\subsection*{Bias: nivel de lectura}
El \texttt{bias} (o nivel cero) mide la señal introducida por la electrónica al
leer el CCD, incluso cuando no llega luz al detector. Se obtiene con el obturador
cerrado y con el menor tiempo de exposición posible, idealmente cercano a cero.
Aunque se espera un nivel aproximadamente constante, la lectura contiene ruido,
por lo que los píxeles presentan pequeñas fluctuaciones alrededor de ese nivel.

\begin{center}
  \begin{tikzpicture}[x=0.62cm,y=0.42cm]
  \draw[->,thick] (0,0) -- (12,0) node[right] {píxeles};
  \draw[->,thick] (0,0) -- (0,4.5) node[above] {cuentas};
  \draw[red,thick] (0.4,2.3) -- (11.4,2.3) node[right] {nivel medio};
  \draw[blue,thick,smooth] plot coordinates {
    (0.5,2.2) (0.9,2.8) (1.3,1.8) (1.7,2.7) (2.1,2.0)
    (2.5,2.6) (2.9,1.9) (3.3,2.8) (3.7,2.1) (4.1,2.6)
    (4.5,1.9) (4.9,2.7) (5.3,2.0) (5.7,2.8) (6.1,2.1)
    (6.5,2.7) (6.9,1.9) (7.3,2.6) (7.7,2.0) (8.1,2.8)
    (8.5,2.1) (8.9,2.6) (9.3,1.9) (9.7,2.7) (10.1,2.0)
    (10.5,2.6) (10.9,2.1)
  };
  \node[align=center] at (6,-0.8) {Fluctuaciones de lectura alrededor de un nivel medio};
\end{tikzpicture}
\end{center}

\subsection*{Dark: corriente oscura}
La imagen \texttt{dark} registra señal no deseada producida principalmente por
la agitación térmica del detector. Su magnitud depende del tiempo de exposición
y de la temperatura. Se toma con el obturador cerrado y con un tiempo igual al
de la imagen científica. Muchos telescopios enfrían el CCD y reducen así la
corriente oscura; aun así, se pueden tomar \texttt{darks} al comienzo o al final
de la noche para comprobar su nivel.

\subsection*{Flat: respuesta píxel a píxel}
Un \texttt{flat} mide cómo responde cada píxel del CCD ante una iluminación
uniforme. Permite corregir diferencias de sensibilidad entre píxeles, polvo u
otras variaciones del sistema óptico y píxeles defectuosos. Se obtiene
apuntando a una fuente uniformemente iluminada, como una pantalla blanca o el
cielo al atardecer. Se debe tomar un conjunto de \texttt{flats} para cada filtro.

\newpage
\section*{2. Calibración y preparación de los datos}

Una vez corregido el nivel de lectura y la señal oscura, se divide por el
\texttt{flat} normalizado. En forma esquemática, para cada píxel:
\[
  I_{\mathrm{limpia}} =
  \frac{I_{\mathrm{cruda}}-I_{\mathrm{bias}}-I_{\mathrm{dark}}}
       {I_{\mathrm{flat, normalizado}}}.
\]
El \texttt{flat} se normaliza para que su nivel medio sea uno; así la división
corrige variaciones relativas sin cambiar arbitrariamente el brillo global de
la imagen.

\begin{center}
  \begin{tikzpicture}[
  box/.style={draw,rounded corners,align=center,minimum width=2.7cm,minimum height=0.9cm},
  arrow/.style={->,thick}
]
  \node[box] (raw) at (0,0) {Imagen\\cruda};
  \node[box] (bias) at (3.6,0) {Restar\\\texttt{bias} y \texttt{dark}};
  \node[box] (flat) at (7.2,0) {Dividir por\\\texttt{flat} normalizado};
  \node[box] (clean) at (10.8,0) {Imagen\\calibrada};
  \draw[arrow] (raw) -- (bias);
  \draw[arrow] (bias) -- (flat);
  \draw[arrow] (flat) -- (clean);
\end{tikzpicture}
\end{center}

\subsection*{Preparación en IRAF}
El ejemplo de las notas utiliza archivos con nombres del tipo
\texttt{m92001.fits}. Primero se descomprime el archivo de datos en la carpeta
de trabajo. Después se inicia IRAF y se cambia a esa carpeta. Los comandos que
se muestran a continuación se ejecutan en la terminal de IRAF:

\begin{verbatim}
rfits fm92** "" junk old+
ls
ds9 &
\end{verbatim}

Se abren varias imágenes en marcos distintos de DS9 para compararlas. Por
ejemplo, se pueden cargar las imágenes \texttt{m92001.fits},
\texttt{m92006.fits} y \texttt{m92015.fits} en marcos separados con
\texttt{display m92001.fits 1}, \texttt{display m92006.fits 2} y
\texttt{display m92015.fits 6}. A continuación se inspeccionan los encabezados:

\begin{verbatim}
imheader m92*.fits l+
imheader m92001.fits l+
\end{verbatim}

El encabezado permite reconocer, entre otros datos, el tipo de imagen
(\texttt{IMAGETYP}, por ejemplo \texttt{'Bias'}), las coordenadas
\texttt{RA} y \texttt{DEC}, el filtro (\texttt{FILTERS}) y la temperatura del
CCD (\texttt{CAMTEMP}). Una línea como
\texttt{m92001.fits [316,509][real]: ave of 25 bias} indica el nombre, el
tamaño de la imagen, el formato numérico y que se trata del promedio de 25
imágenes \texttt{bias}.

\newpage
\section*{3. Inspección de imágenes y perfiles con \texttt{implot}}

Se visualizan las imágenes en DS9 y se describen sus rasgos generales. Para
examinar los valores de los píxeles a lo largo de una fila o columna se usa
\texttt{implot}. Por ejemplo:

\begin{verbatim}
implot m92001
\end{verbatim}

En el ejemplo se inspecciona la línea 255. El perfil puede comenzar con un
aumento, disminuir después y mostrar fluctuaciones alrededor de un nivel de
aproximadamente 500 cuentas. Promediar varias líneas ayuda a distinguir la
señal sistemática del ruido. En la ventana de \texttt{implot}, la combinación
\texttt{Shift} + \texttt{:} abre la entrada de comandos; se puede solicitar,
por ejemplo, el promedio de las líneas 10 a 500 con \texttt{: l 10 500}.

\begin{center}
  \begin{tikzpicture}[x=0.68cm,y=0.32cm]
  \draw[->,thick] (0,0) -- (12,0) node[right] {columna};
  \draw[->,thick] (0,0) -- (0,7) node[above] {cuentas};
  \draw[gray,dashed] (8.6,0) -- (8.6,6.4);
  \draw[blue,thick,smooth] plot coordinates {
    (0.3,3.0) (0.8,3.2) (1.3,3.1) (1.8,3.4) (2.3,3.2)
    (2.8,3.1) (3.3,3.3) (3.8,3.2) (4.3,3.4) (4.8,3.1)
    (5.3,3.3) (5.8,3.2) (6.3,3.1) (6.8,3.4) (7.3,3.2)
    (7.8,3.3) (8.3,3.2) (8.7,4.0) (9.0,6.2) (9.3,3.5)
    (9.8,3.2) (10.3,3.4) (10.8,3.1) (11.3,3.3)
  };
  \node[align=center] at (4.4,5.3) {Área útil};
  \node[align=center] at (10.1,5.3) {pico en el\\\texttt{overscan}};
  \node[below] at (8.6,-0.15) {borde};
\end{tikzpicture}
\end{center}

Al inspeccionar otra imagen, como \texttt{m92006}, se repite el procedimiento:
\texttt{display m92006}, \texttt{implot m92006} y el promedio de líneas que
corresponda. En DS9 puede aparecer un rayo vertical que coincide con un pico
del perfil de \texttt{implot}; en el ejemplo, se encuentra después de la
columna 300. Comparar ambas visualizaciones ayuda a localizar estructuras del
detector que conviene excluir de la región científica.

\subsection*{La región overscan}
En algunos detectores se leen filas o columnas adicionales, aunque no formen
parte del área fotosensible. Esta región ficticia, llamada \texttt{overscan},
permite seguir el nivel electrónico de lectura (el \texttt{bias} del
amplificador). En el perfil aparece como una subida o caída marcada al llegar
al borde de la imagen. Los \texttt{overscans} también pueden observarse en los
\texttt{flats}; en las imágenes del conjunto de ejemplo se señalan en
\texttt{m92007} y \texttt{m92010}.

\newpage
\section*{4. Identificación y recorte del overscan}

Para tener más información sobre el nivel de lectura, el detector puede
incorporar una zona de lectura adicional en uno o ambos bordes. Como no
corresponde a la escena observada, esa zona debe identificarse y excluirse antes
del análisis científico. En estas notas el exceso aparece en columnas, por lo
que se emplea la tarea \texttt{colbias}; si el exceso estuviera en filas, se
consideraría una corrección de tipo \texttt{linebias}.

\begin{center}
  \begin{tikzpicture}[x=0.65cm,y=0.65cm]
  \draw[thick] (0,0) rectangle (10,5.2);
  \fill[blue!10] (0,0) rectangle (8.6,5.2);
  \fill[red!18] (8.6,0) rectangle (10,5.2);
  \draw[thick,dashed] (8.6,0) -- (8.6,5.2);
  \node[align=center] at (4.3,2.6) {Región científica};
  \node[align=center] at (9.3,2.6) {\texttt{overscan}};
  \draw[<->,thick] (0,-0.55) -- (8.6,-0.55)
    node[midway,below] {\texttt{TRIM}: conservar};
  \draw[->,thick] (9.3,5.3) -- (9.3,6.15)
    node[above,align=center] {Estimar el nivel\\con \texttt{BIAS}};
\end{tikzpicture}
\end{center}

La secuencia de consulta y configuración de IRAF es:

\begin{verbatim}
help colbias
noao
imred
bias
unlearn colbias
lpar colbias
help colbias
\end{verbatim}

\texttt{help colbias} muestra la ayuda de la tarea; \texttt{noao},
\texttt{imred} y \texttt{bias} permiten entrar en los paquetes pertinentes.
\texttt{unlearn colbias} restablece los parámetros, y \texttt{lpar colbias}
los lista para su revisión. La sección de ejemplos de la ayuda muestra cómo
asignar la región que se usará para el nivel de \texttt{bias}; por ejemplo, la
sintaxis puede escribirse como:

\begin{verbatim}
cl> colbias.bias = "[301:832]"
\end{verbatim}

La región se expresa como \texttt{[x1:x2,y1:y2]}, donde el primer intervalo
corresponde a columnas y el segundo a filas. Conviene confirmar la orientación
de los ejes y los límites con la imagen concreta antes de aplicar la tarea.

\newpage
\section*{5. Selección de las regiones \texttt{BIAS} y \texttt{TRIM}}

Los límites no se deben elegir a ciegas: se inspecciona el perfil con
\texttt{implot} y se hace zoom sobre los bordes de la imagen. En la ventana de
\texttt{implot}, \texttt{e} amplía la región seleccionada, \texttt{r} permite
volver a la vista anterior y \texttt{q} cierra la tarea. La combinación
\texttt{Shift} + \texttt{?} muestra las funciones disponibles.

Para elegir las filas o líneas que se promediarán, se puede examinar una línea
con \texttt{c 225} y promediar un intervalo con \texttt{c 10 300}. El perfil
debe mostrar dónde termina el área útil y dónde comienza el nivel de lectura
del borde. En el ejemplo, la región de \texttt{bias} ocupa las columnas 330 a
350 y las filas 2 a 510:

\[
  \texttt{BIAS} = [330:350,\,2:510].
\]

El área científica se conserva mediante la región \texttt{TRIM}. Para el
ejemplo, se retienen las columnas 1 a 318 y las filas 2 a 510:

\[
  \texttt{TRIM} = [1:318,\,2:510].
\]

Los valores anteriores son ilustrativos y deben verificarse en cada imagen:
se amplía el borde izquierdo para localizar el comienzo del área útil y el
borde derecho para ubicar el final de la región científica y el
\texttt{overscan}. Una vez delimitadas ambas regiones, se configura
\texttt{colbias} para estimar y sustraer el nivel de lectura, y se recorta la
imagen a \texttt{TRIM} para continuar el análisis solo con los píxeles de la
escena.


\newpage
\section*{6. Aplicación de \texttt{colbias} por lotes}

Una vez definidos los intervalos \texttt{BIAS} y \texttt{TRIM}, se guardan en
los parámetros de la tarea \texttt{colbias}. Desde IRAF se abre el editor de
parámetros con \texttt{epar colbias} y se revisan las opciones:

\begin{itemize}
  \item \texttt{function = chebyshev}: ajusta un polinomio de Chebyshev al
  nivel de lectura. El rechazo bajo y alto de tres desviaciones estándar
  descarta valores extremos; con dos desviaciones, la desigualdad de
  Chebyshev garantiza al menos el 75\% de los datos para una distribución
  general.
  \item \texttt{interactive = no} y \texttt{median = no}: no se solicita un
  ajuste interactivo ni se usa la mediana.
  \item \texttt{order = 1}: se usa un ajuste de primer orden.
  \item \texttt{low\_reject = 3.0} y \texttt{high\_reject = 3.0}: fijan los
  umbrales de rechazo en ambos lados.
\end{itemize}

Al terminar, se elige \texttt{Save and Quit}. Para procesar un conjunto de
archivos se preparan dos listas: una con los nombres de entrada y otra con los
nombres de salida. La segunda contiene el prefijo \texttt{Trm} para distinguir
las imágenes recortadas y corregidas.

\begin{verbatim}
ls m92*.fits
ls m92*.fits > input
\end{verbatim}

Se copia la lista \texttt{input}, se cambia el nombre de la copia a
\texttt{output} y se añade \texttt{Trm} al inicio de cada nombre. Después se
indica a \texttt{colbias} que lea ambas listas y se ejecuta la tarea:

\begin{verbatim}
epar colbias
input  = @input
output = @output
\end{verbatim}

\begin{center}
  \begin{tikzpicture}[
  box/.style={draw,rounded corners,align=center,minimum width=2.5cm,minimum height=0.85cm},
  arrow/.style={->,thick}
]
  \node[box] (input) at (0,0) {Lista \texttt{input}\\imágenes originales};
  \node[box] (task) at (4,0) {\texttt{colbias}\\\texttt{BIAS} y \texttt{TRIM}};
  \node[box] (output) at (8,0) {Lista \texttt{output}\\prefijo \texttt{Trm}};
  \node[box] (images) at (12,0) {Archivos\\corregidos};
  \draw[arrow] (input) -- (task);
  \draw[arrow] (task) -- (output);
  \draw[arrow] (output) -- (images);
\end{tikzpicture}
\end{center}

El resultado es un conjunto de archivos \texttt{Trm92*.fits}. Se comprueba que
se generaron y se inspecciona una imagen para confirmar que el recorte eliminó
la franja del \texttt{overscan} sin quitar la región útil.

\newpage
\section*{7. Verificación y clasificación de las imágenes}

Se comparan en DS9 una imagen original y su resultado de \texttt{colbias}:

\begin{verbatim}
display m92010.fits 1
display Trm92010.fits 2
\end{verbatim}

En el ejemplo se comprueba que se eliminó la zona negra del borde. También se
observa que las cuentas pueden seguir siendo altas: \texttt{colbias} corrige
el nivel del borde y recorta la imagen, pero no sustituye la corrección de
\texttt{flat}. Conviene comprobar además que no se haya perdido una región
iluminada del lado izquierdo.

La clasificación se confirma con los encabezados FITS. En el conjunto de
práctica se identifican los siguientes tipos de imagen:

\begin{center}
\begin{tabular}{ll}
\hline
Archivo & Tipo y filtro\\
\hline
\texttt{m92001} & \texttt{bias}\\
\texttt{m92006} & \texttt{flat}, filtro B\\
\texttt{m92007} & \texttt{flat}, filtro V\\
\texttt{m92010}, \texttt{m92011} & Imágenes científicas, filtro V\\
\texttt{m92014}, \texttt{m92015} & Imágenes científicas, filtro B\\
\hline
\end{tabular}
\end{center}

Las imágenes \texttt{m92010} y \texttt{m92011} se comparan con el
\texttt{flat} V; las imágenes \texttt{m92014} y \texttt{m92015}, con el
\texttt{flat} B. Mantener separadas las bandas evita aplicar una corrección de
sensibilidad que no corresponde al filtro usado en la observación.

\newpage
\section*{8. Normalización de los \texttt{flats}}

En este conjunto no hay imágenes \texttt{dark}: las observaciones se hicieron
con un CCD refrigerado y, según las notas, se omite esa corrección. El siguiente
paso es normalizar cada \texttt{flat} para poder dividir las imágenes
científicas por él.

Un \texttt{flat} ideal tendría la misma cuenta en todos sus píxeles. En la
práctica, los píxeles responden de manera algo diferente y el nivel de cuentas
varía. Se estima un valor representativo con la moda (el valor más frecuente)
y se divide el \texttt{flat} por esa constante.

\begin{center}
  \begin{tikzpicture}[x=0.48cm,y=0.48cm]
  \node[font=\bfseries] at (2,4.7) {Flat ideal};
  \node[font=\bfseries] at (9,4.7) {Flat observado};
  \draw (0,0) grid (4,4);
  \draw (7,0) grid (11,4);
  \foreach \x in {0,...,3} {
    \foreach \y in {0,...,3} {
      \node at (\x+0.5,\y+0.5) {5};
    }
  }
  \node at (7.5,3.5) {5}; \node at (8.5,3.5) {4};
  \node at (9.5,3.5) {6}; \node at (10.5,3.5) {5};
  \node at (7.5,2.5) {4}; \node at (8.5,2.5) {5};
  \node at (9.5,2.5) {5}; \node at (10.5,2.5) {6};
  \node at (7.5,1.5) {6}; \node at (8.5,1.5) {5};
  \node at (9.5,1.5) {4}; \node at (10.5,1.5) {5};
  \node at (7.5,0.5) {5}; \node at (8.5,0.5) {6};
  \node at (9.5,0.5) {5}; \node at (10.5,0.5) {4};
  \draw[->,thick] (4.3,2) -- (6.7,2);
  \node[align=center] at (5.5,-0.8) {Variación entre píxeles};
\end{tikzpicture}
\end{center}

La tarea \texttt{imstat} resume estadísticas de las imágenes. Se puede
consultar su ayuda y solicitar, por ejemplo, el nombre, la media y la moda:

\begin{verbatim}
imstat BTrm92006.fits
help imstat
imstat BTrm92006.fits fields="image,mean,mode"
imstat BTrm92006.fits,BTrm92007.fits fields="image,mode"
\end{verbatim}

Para las imágenes de ejemplo, las notas usan 1312 cuentas para el filtro B y
1467 para el filtro V. Se divide cada \texttt{flat} por su propio valor y se
crean los archivos \texttt{NFlatB} y \texttt{NFlatV}:

\begin{verbatim}
imarith BTrm92006.fits / 1312 NFlatB.fits
imarith BTrm92007.fits / 1467 NFlatV.fits
\end{verbatim}

Estas constantes son específicas de los datos de práctica; antes de usarlas en
otras imágenes se debe volver a calcular la moda. La normalización cambia la
escala de las cuentas, aunque la apariencia general del \texttt{flat} pueda
parecer similar.

\newpage
\section*{9. Corrección con el \texttt{flat} correspondiente}

Se divide cada imagen científica por el \texttt{flat} normalizado del mismo
filtro. Por tanto, se usa \texttt{NFlatV} solo para las observaciones V y
\texttt{NFlatB} solo para las observaciones B. El prefijo \texttt{F} en los
archivos de salida señala que se aplicó la corrección de \texttt{flat}.

\begin{center}
  \begin{tikzpicture}[
  box/.style={draw,rounded corners,align=center,minimum width=2.8cm,minimum height=0.8cm},
  arrow/.style={->,thick}
]
  \node[font=\bfseries] at (0,1.8) {Filtro V};
  \node[box] (vscience) at (2.3,1.8) {Imágenes V\\\texttt{BTrm92010}\\\texttt{BTrm92011}};
  \node[box] (vflat) at (6.2,1.8) {Dividir por\\\texttt{NFlatV}};
  \node[box] (vout) at (10.2,1.8) {Salida\\\texttt{FBTrm...}};
  \draw[arrow] (vscience) -- (vflat);
  \draw[arrow] (vflat) -- (vout);
  \node[font=\bfseries] at (0,-0.2) {Filtro B};
  \node[box] (bscience) at (2.3,-0.2) {Imágenes B\\\texttt{BTrm92014}\\\texttt{BTrm92015}};
  \node[box] (bflat) at (6.2,-0.2) {Dividir por\\\texttt{NFlatB}};
  \node[box] (bout) at (10.2,-0.2) {Salida\\\texttt{FBTrm...}};
  \draw[arrow] (bscience) -- (bflat);
  \draw[arrow] (bflat) -- (bout);
\end{tikzpicture}
\end{center}

La operación \texttt{imarith} se aplica a cada imagen; algunos ejemplos son:

\begin{verbatim}
imarith BTrm92010.fits / NFlatV.fits FBTrm92010.fits
imarith BTrm92011.fits / NFlatV.fits FBTrm92011.fits
imarith BTrm92014.fits / NFlatB.fits FBTrm92014.fits
imarith BTrm92015.fits / NFlatB.fits FBTrm92015.fits
\end{verbatim}

Se inspecciona el resultado comparando la imagen antes y después de dividir por
el \texttt{flat}, por ejemplo, \texttt{display m92010.fits 1} y
\texttt{display FBTrm92010.fits 2}. También se puede visualizar el propio
\texttt{flat} normalizado junto con el original para comprobar que la
normalización conserva su patrón espacial y cambia la escala de los valores.

\newpage
\section*{10. Inicio de la fotometría y masa de aire}

Con las imágenes corregidas se inicia la etapa de fotometría. Primero se
revisan los encabezados y se visualizan imágenes de cada filtro:

\begin{verbatim}
imheader FBTrm92*.fits l-
ds9 &
display FBTrm92010.fits 1
display FBTrm92014.fits 2
\end{verbatim}

La masa de aire describe la cantidad relativa de atmósfera atravesada por la
luz desde el objeto hasta el telescopio. Es mínima cerca del cenit y aumenta
hacia el horizonte; debe tenerse en cuenta al preparar la fotometría y
caracterizar la función de dispersión puntual (PSF).

\begin{center}
  \begin{tikzpicture}[x=0.85cm,y=0.62cm]
  \fill[blue!8] (-0.2,0.5) rectangle (10,1.1);
  \draw[thick] (-0.2,0.5) -- (10,0.5);
  \node[below] at (4.8,0.5) {superficie terrestre};
  \node[draw,circle,fill=black,inner sep=1.5pt,label=below:{telescopio}] (obs) at (2,0.5) {};
  \draw[->,very thick,blue] (2,0.5) -- (2,5.2);
  \draw[->,very thick,red] (2,0.5) -- (8.6,3.9);
  \node[star,star points=5,fill=yellow!80!orange,draw,minimum size=0.45cm] at (2,5.2) {};
  \node[star,star points=5,fill=yellow!80!orange,draw,minimum size=0.45cm] at (8.6,3.9) {};
  \draw[dashed] (2,0.5) -- (8.6,0.5);
  \draw (2,1.55) arc[start angle=90,end angle=63,radius=1.05];
  \node at (2.55,1.8) {$z$};
  \node[align=center,blue] at (1.15,3.1) {cerca del\\cenit: menor\\masa de aire};
  \node[align=center,red] at (7.1,2.2) {trayecto oblicuo\\más largo};
\end{tikzpicture}
\end{center}

Para completar los encabezados se consulta la tarea \texttt{setairmass} y se
entra en el paquete de utilidades astronómicas de IRAF. Las notas también
proponen guardar el encabezado de una imagen en un archivo de texto para
inspeccionarlo:

\begin{verbatim}
help setairmass
noao
astutil
imheader FBTrm92010.fits > header
\end{verbatim}

El archivo \texttt{header} permite revisar la información disponible antes de
asignar o verificar la masa de aire de las imágenes.

\end{document}
